\section{Before the Fathers: Israel's Writings and the Philosophers}\label{sec:before-christ}

The Fathers learned the angels from the Scriptures, and they read the
Scriptures within an inheritance wider than the canon. The Greek Bible of
the Apostles names the angels where the Hebrew names the sons of Israel or
the sons of God. The writers of Israel in the last centuries before Christ
told of the Watchers who fell, of the seven who stand before God, of the
angels of the nations and the angel who intercedes for Israel. Philo of
Alexandria set Moses' angels beside the daemons of the philosophers, and
the poets and philosophers of Greece taught of spirits between the gods
and men. The Fathers knew all of this. They received what agreed with the
Scriptures and the rule of faith, corrected what was mixed with error, and
rejected what contradicted the faith, and in each case the judgment was
made by Scripture as the Church received it and by the faith she confesses:
that the angels are creatures of the one God, made good, free, and sent to
serve the salvation of men, and that the demons are angels who fell by
their own will.

\sectionguard
\subsection{The Angels of the Septuagint}\label{sec:bc-lxx}

The Greek translation of the Law and the Prophets was the Bible of the
Apostles' mission to the nations and of the Greek Fathers, and its Psalter
passed into the Latin Church through the Vulgate. Augustine holds that its
departures from the Hebrew may themselves be the Holy Spirit's work: ``the
Septuagint translators are justly believed to have received the Spirit of
prophecy; so that, if they made any alterations under His authority, and
did not adhere to a strict translation, we could not doubt that this was
divinely dictated'' (\emph{De civitate Dei} XV.23). Several of those
departures concern the angels, and on them the Fathers built.

\paragraph{The angels of the nations (Deut 32:8--9).} In the Song of Moses
the Greek reads: ``When the Most High divided the nations, when he
separated the sons of Adam, he set the bounds of the nations according to
the number of the angels of God. And his people Jacob became the portion of
the Lord, Israel was the line of his inheritance'' (Deut 32:8--9,
Septuagint, Brenton). The Douay, with the Vulgate and the Hebrew, reads
that God ``appointed the bounds of people according to the number of the
children of Israel.'' Clement of Rome quotes the Greek form (\emph{1 Clem.}\ 29), and Justin,
arguing with Trypho from the Hebrew, notes: ``The Seventy have translated it,
`He set the bounds of the nations according to the number of the angels of
God'\,'' (\emph{Dial.}\ 131). Clement of Alexandria teaches that ``by an
ancient and divine order the angels are distributed among the nations. But
the glory of those who believe is `the Lord's portion'\,'' (\emph{Strom.}\
VII.2). Origen reads the verse with the tower of Babel: the nations ``are
called a part of the angels,'' handed over at the dispersion to the angels
who led them to their lands (\emph{De principiis} I.5.2; \emph{C.\ Cels.}\
V.29--30; see \ref{sec:fathers-before}). Basil joins the Song to Daniel's
princes of Persia and of the Greeks: some angels are set over whole
nations, one angel rules and another is subject, and all are angels by
nature (\emph{Adversus Eunomium} III.1, PG~29).

Eusebius draws the Christological consequence. The nations ``were
distributed among the invisible guardians of the nations, that is the
angels, by the decision of the Most High God, and His secret counsel unknown
to us. Whereas to One beyond comparison with them, the Head and King of the
Universe, I mean to Christ Himself, as being the Only-begotten Son, was
handed over that part of humanity denominated Jacob and Israel''
(\emph{Demonstratio evangelica} IV.7). Irenaeus and Hilary read the same
verse of the Gospel's going out to the nations. Paul at Athens preached one
God who made one race of men, ``as also Moses declared,'' and ``that people
which believes in God is not now under the power of angels, but under the
Lord's [rule]'' (Irenaeus, \emph{Adv.\ haer.}\ III.12.9). Hilary, on the
promise ``I will give thee the Gentiles for thy inheritance'' (Ps 2:8),
says that the nations baptized and taught are \latin{non iam secundum
diuinam illam Moysi cantionem angelorum dominatui deditae neque secundum
eorundem numerum diuisae, sed in dominicam familiam susceptae}: no longer
given over to the rule of angels according to Moses' divine song, nor
divided by their number, but taken into the Lord's household (\emph{Tract.\
in Ps.}\ 2.31).

Dionysius gives the verse its place in the hierarchy. The lowest order has
charge of men, and so ``the Word of God has assigned our Hierarchy to
Angels, by naming Michael as Ruler of the Jewish people, and others over
other nations. For the Most High established borders of nations according
to number of Angels of God'' (\emph{CH} 9.2; Dan 10:13, 21). He at once
excludes a reading the verse might seem to allow. If the nations went
after false gods, ``we ought not to throw the blame of the other nations
wandering after those which are no gods upon the direct guidance of the
Angels, but \dots\ they themselves, by their own declension, fell away''
(\emph{CH} 9.3); and ``never must we think that the Godhead is leader of Jews
by lot, and that Angels, independently, or as of equal rank, or in
opposition, or that certain other gods, preside over the other nations.''
There is one Providence over all, and ``all the Angels who preside over
each nation, elevate, as far as possible, those who follow them with a
willing mind, to It as their proper Head'' (\emph{CH} 9.4; see
\ref{sec:ch9}). Aquinas receives the teaching through the order of
Principalities, ``who, as presiding over the government of peoples and
kingdoms (which occupies the first and principal place in the Divine
ministrations), are the first in the execution thereof'' (\ST{I q.~108
a.~6}), and he holds with Gregory that ``the prince of the kingdom of
Persia was a good angel appointed to the guardianship of that kingdom''
(\ST{I q.~113 a.~8}; see \ref{sec:q108}, \ref{sec:q113}). The Catechism
teaches the same of the nations after Babel, ``each entrusted by divine
providence to the guardianship of angels,'' and cites for it the Septuagint
of Deut 32:8 (CCC~57).

\paragraph{The angels of God before the Lord (Job 1:6; 2:1; 38:7).} Where
the Hebrew and the Douay say that ``the sons of God came to stand before the
Lord, Satan also was present among them,'' the Greek reads: ``behold, the
angels of God came to stand before the Lord, and the devil came with
them'' (Job 1:6, Septuagint; so also 2:1). Aquinas reads the Latin text in
the Greek's sense: Satan ``is not described as having assisted, but as
present among the assistants,'' having kept ``a nature like to the angels''
(\ST{I q.~112 a.~3 ad~3}, with Gregory). At the founding of the world the
Douay hears that ``the morning stars praised me together, and all the sons
of God made a joyful melody'' (Job 38:7); the Greek reads, ``When the stars
were made, all my angels praised me with a loud voice.'' The Greek verse
became the Fathers' chief argument for the time of the angels' creation,
and they drew opposite conclusions from it. Augustine: since the angels
praised God when the stars were made on the fourth day, ``The angels
therefore existed before the stars,'' and being among the works of the six
days they are ``that light which was called `Day'\,'' (\emph{De civ.\ Dei}
XI.9). Cassian's Abbot Serenus: ``Those then who were present at the
creation of the stars, are most clearly proved to have been created before
that `beginning' in which it is said that heaven and earth were made''
(\emph{Conl.}\ VIII.7; see \ref{sec:fathers-latin}). Aquinas records both
opinions: the Greek Fathers ``all \dots\ hold the
creation of the angels to have taken place previously to that of the
corporeal world,'' while the more probable opinion holds that they were
created with it, for ``the angels are part of the universe''; the contrary
``is not to be deemed erroneous'' (\ST{I q.~61 a.~3}; see \ref{sec:q61}).
The Fourth Lateran Council confesses that God \latin{simul ab initio temporis
utramque de nihilo condidit creaturam, spiritualem et corporalem, angelicam
videlicet et mundanam} (DS~800).

At Gen 6:2 the Greek manuscripts divide: Brenton's text has ``the sons of
God,'' and he notes that the Alexandrine manuscript reads ``angels of God.''
Augustine knew the variation: ``The Septuagint indeed calls them both angels
of God and sons of God, though all the copies do not show this, some having
only the name `sons of God'\,'' (\emph{De civ.\ Dei} XV.23). On this reading
the Watchers' history turns (\ref{sec:bc-enoch}).

\paragraph{The worship of the angels given to the Son.} The Greek Song of
Moses ends, ``let all the angels of God worship him'' (Deut 32:43,
Septuagint), and the Greek Psalter has ``worship him, all ye his angels'' (Ps
96[97]:7); Hebrews takes the words of the Son's entry into the world: ``And
let all the angels of God adore him'' (Heb 1:6). The Vulgate Psalter,
following the Greek, keeps the angels where the Hebrew speaks of God or the
mighty: man was made ``a little less than the angels'' (Ps 8:6; Heb 2:7); in
the desert ``Man ate the bread of angels'' (Ps 77[78]:25; Wis 16:20); ``I will
sing praise to thee in the sight of the angels'' (Ps 137[138]:1). From this
last verse Benedict gives his monks their rule of psalmody: \latin{In
conspectu angelorum psallam tibi. Ergo consideremus qualiter oporteat in
conspectu Divinitatis et angelorum eius esse}, and so stand to sing that the
mind agrees with the voice (\emph{Regula} 19).

\paragraph{The Law given through angels (Deut 33:2).} Where the Douay reads
that the Lord came from Sinai with ``In his right hand a fiery law,'' the
Greek has ``on his right hand were his angels with him'' (Deut 33:2,
Septuagint). The New Testament teaches that Israel ``received the law by the
disposition of angels'' (Acts 7:53), ``ordained by angels in the hand of a
mediator'' (Gal 3:19), ``the word spoken by angels'' (Heb 2:2). On these texts
Aquinas determines that ``God made the Law by His own authority, but He
promulgated it through the angels,'' while the perfect law was given by ``the
incarnate God immediately'' (\ST{I-II q.~98 a.~3} and ad~3).

\paragraph{The Angel of great counsel (Isa 9:6; 63:9).} Where the Douay
names the Child ``Wonderful, Counsellor, God the Mighty,'' the Greek reads,
``his name is called the Messenger of great counsel'' (Isa 9:6, Septuagint;
\gk{megal\=es boul\=es angelos}). Justin counts ``the Angel of great
counsel'' among the titles of Christ (\emph{Dial.}\ 126), and Origen answers
Celsus that the prophecy calls Jesus ``not simply an angel, but the `Angel of
the great counsel:' for He announced to men the great counsel of the God and
Father of all things'' (\emph{C.\ Cels.}\ V.53). At Isa 63:9, where the Douay
reads that ``the angel of his presence saved them,'' the Greek has ``not an
ambassador, nor a messenger, but himself saved them,'' and Irenaeus reads it
of the Incarnation: not a mere man ``nor [one] without flesh---for the angels
are without flesh'' was to save us, but ``the Lord Himself'' (\emph{Adv.\
haer.}\ III.20.4).

\paragraph{The furnace and the watcher.} The Greek portion of Daniel, which
the Vulgate carries, adds to the story of the three young men that ``the
angel of the Lord went down with Azarias and his companions into the
furnace,'' and the king blesses the God ``who hath sent his angel, and
delivered his servants'' (Dan 3:49, 95). In their canticle the young men
call the angels to praise next after the heavens: ``O ye angels of the
Lord, bless the Lord: praise and exalt him above all for ever'' (Dan 3:58).
Augustine proves from it that the angels are God's works: among the works
called to bless the Lord, ``the angels are named'' (\emph{De civ.\ Dei}
XI.9).

In Nabuchodonosor's dream the Aramaic chapters of Daniel give the angels a
name of their own: ``behold a watcher, and a holy one came down from
heaven''; the sentence on the king is ``the decree by the sentence of the
watchers'' (Dan 4:10, 14 [4:13, 17]). Jerome explains the word. For \latin{vigil} Theodotion kept the
Aramaic itself, \emph{ir}, which \latin{significat autem angelos quod semper
vigilent; et ad Dei imperium sint parati. Unde et nos crebris
pernoctationibus imitamur angelorum officia}: it signifies the angels,
because they are ever wakeful and ready at God's command; and so we too, by
frequent night vigils, imitate the office of the angels (\emph{In Dan.}\
4:10). Watcher is the name the Enochic books give to the angels, holy and
fallen alike.

\sectionguard
\subsection{The Books of Enoch and the Watchers}\label{sec:bc-enoch}

The Book of Enoch gathers under the name of the patriarch who ``walked with
God'' (Gen 5:24) the vision of the heavenly court, the fall of the Watchers
and their judgment, Enoch's journeys through heaven and Sheol, and the
parables of the judgment to come.\footnote{English of 1~Enoch from R. H.
Charles, \emph{The Book of Enoch} (London: SPCK, 1917), with the
translator's critical signs omitted.} The Epistle of Jude quotes it:
``Now of these Enoch also, the seventh from Adam, prophesied, saying:
Behold, the Lord cometh with thousands of his saints: To execute judgment
upon all'' (Jude 14--15). The prophecy stands at the head of the book:
``And behold! He cometh with ten thousands of His holy ones To execute
judgement upon all, And to destroy all the ungodly'' (1~En 1:9). Jude's
``angels who kept not their principality but forsook their own habitation,''
reserved ``under darkness in everlasting chains, unto the judgment of the
great day'' (Jude 6), and Peter's angels ``drawn down by infernal ropes to
the lower hell \dots\ to be reserved unto judgment'' (2~Pet 2:4) speak in
the book's terms of darkness, bonds, and judgment; Augustine refers them to
the first fall of the devil and his angels (\emph{De civ.\ Dei} XV.23). The
Council of Trent received Jude among
the canonical books (DS~1503); the Book of Enoch the Church did not receive,
and the Fathers weighed it each by his own judgment.

\paragraph{The fall of the Watchers (1~En 6--8).} The Book of the Watchers
tells the history of Gen 6:1--4 as a history of angels. ``And the angels,
the children of the heaven, saw and lusted after them, and said to one
another: `Come, let us choose us wives from among the children of men and
beget us children'\,'' (1~En 6:2). Their leader Semjâzâ fears to bear the
penalty alone, and they bind themselves by an oath: ``they were in all two
hundred; who descended in the days of Jared on the summit of Mount Hermon,
and they called it Mount Hermon, because they had sworn and bound themselves
by mutual imprecations upon it'' (6:6). They took wives and taught them
charms, enchantments, and the cutting of roots, and the women bore giants
who devoured the labour of men and at last men themselves (7:1--5). Azâzêl
``taught men to make swords, and knives, and shields, and breastplates, and
made known to them the metals of the earth and the art of working them, and
bracelets, and ornaments, and the use of antimony, and the beautifying of
the eyelids''; others taught astrology and the signs of sun and moon
(8:1--3).

\paragraph{The judgment of the Watchers (1~En 9--16).} Michael, Uriel,
Raphael, and Gabriel look down and see the blood shed on the earth, and
they carry the cry of men before God: ``to you, the holy ones of heaven, the
souls of men make their suit, saying, `Bring our cause before the Most
High'\,'' (9:1--3). God answers by his archangels. Raphael is to ``Bind
Azâzêl hand and foot, and cast him into the darkness,'' where he will lie
until ``the day of the great judgement'' when ``he shall be cast into the
fire'' (10:4--6); Gabriel is sent against the giants; Michael is to bind
Semjâzâ and his fellows ``fast for seventy generations in the valleys of the
earth, till the day of their judgement'' (10:11--12). Enoch is sent to ``the
Watchers of the heaven who have left the high heaven, the holy eternal
place, and have defiled themselves with women'' (12:4). They beg him to
intercede, and God's answer states the order they broke: ``You should
intercede for men, and not men for you'' (15:2). ``But you were formerly
spiritual, living the eternal life, and immortal for all generations of the
world. And therefore I have not appointed wives for you; for as for the
spiritual ones of the heaven, in heaven is their dwelling'' (15:6--7). Of
the giants' spirits the book says, ``the giants, who are produced from the
spirits and flesh, shall be called evil spirits upon the earth, and on the
earth shall be their dwelling'' (15:8). The fallen ``have been in heaven,
but all the mysteries had not yet been revealed to'' them, and what they
knew they ``made known to the women'' (16:3). Uriel shows Enoch the place
where the angels who joined themselves to women shall stand, whose ``spirits
assuming many different forms are defiling mankind and shall lead them
astray into sacrificing to demons as gods'' (19:1).

\paragraph{The holy Watchers.} The same book shows the unfallen angels in
their service. Before the throne of the Great Glory ``ten thousand times ten
thousand (stood) before Him, yet He needed no counselor. And the most holy
ones who were nigh to Him did not leave by night nor depart from Him''
(14:22--23; Dan 7:10). In the Parables the court sings the song of Isaias:
``Those who sleep not bless Thee: they stand before Thy glory and bless,
praise, and extol, saying: `Holy, holy, holy, is the Lord of Spirits: He
filleth the earth with spirits'\,'' (39:12; Isa 6:3). Around the throne are
``Seraphin, Cherubic, and Ophannin: And these are they who sleep not And
guard the throne of His glory'' (71:7), and with them God summons ``all the
angels of power, and all the angels of principalities'' (61:10), names that
stand beside the Apostle's (Col 1:16). Four presences surround the Lord, one
of whom intercedes for those on earth and another fends off ``the Satans''
from accusing them: ``This first is Michael, the merciful and long-suffering:
and the second, who is set over all the diseases and all the wounds of the
children of men, is Raphael: and the third, who is set over all the powers,
is Gabriel: and the fourth, who is set over the repentance unto hope of those
who inherit eternal life, is named Phanuel'' (40:6--9). The righteous are
promised that God ``will appoint guardians from amongst the holy angels To
guard them as the apple of an eye'' (100:5).

\paragraph{The Fathers who received the Watchers.} The second-century
apologists tell the fall of the angels in the book's terms, without naming
it. Justin teaches that God ``committed the care of men and of all things
under heaven to angels whom He appointed over them. But the angels
transgressed this appointment, and were captivated by love of women, and
begat children who are those that are called demons'' (\emph{2~Apol.}\ 5).
Athenagoras, who says ``we say nothing without witnesses, but state the
things which have been declared by the prophets,'' tells of angels who ``fell
into impure love of virgins'' and begot the giants (\emph{Legatio} 24), and
distinguishes, as Enoch does (1~En 15:8--16:1), the fallen angels, who
``haunt the air and the earth,'' from ``the souls of the giants, which are
the demons who wander about the world'' (\emph{Leg.}\ 25).
Irenaeus knows Enoch's embassy to the Watchers (1~En 12--16): Enoch
``discharged the office of God's legate to the angels although he was a man
\dots\ because the angels when they had transgressed fell to the earth for
judgment'' (\emph{Adv.\ haer.}\ IV.16.2; cf.\ IV.36.4). Clement of Alexandria gives
the Enochic account of the secrets: ``the angels who had obtained the superior rank, having sunk into
pleasures, told to the women the secrets which had come to their knowledge;
while the rest of the angels concealed them, or rather, kept them against
the coming of the Lord'' (\emph{Strom.}\ V.1; 1~En 16:3), and in his notes on
Jude he takes the Apostle's words as confirming the prophecy
(\emph{Adumbr.\ in Iud.}).

Tertullian names the book and defends it. The ornaments of women were
first given by ``those angels, to wit, who rushed from heaven on the
daughters of men,'' who ``had laid bare the operations of metallurgy, and
had divulged the natural properties of herbs, and had promulgated the powers
of enchantments, and had traced out every curious art, even to the
interpretation of the stars'' (\emph{De cultu feminarum} I.2; 1~En 8:1--3).
To those who refuse ``the Scripture of Enoch'' because the Jewish canon does
not admit it, he answers that ``since Enoch in the same Scripture has
preached likewise concerning the Lord, nothing at all must be rejected by us
which pertains to us,'' and that ``Enoch possesses a testimony in the
Apostle Jude'' (\emph{De cultu fem.}\ I.3). Against idolatry he quotes it
twice: ``Enoch had preceded, predicting that `the demons, and the spirits
of the angelic apostates, would turn into idolatry all the elements \dots\
that they might be consecrated as God, in opposition to God'\,''; and again,
``Ye who serve stones, and ye who make images of gold, and silver, and wood,
and stones and clay, and serve phantoms, and demons, and spirits in fanes
\dots\ shall find no help from them'' (\emph{De idololatria} 4; 1~En 19:1;
99:7). The astrologers are condemned with ``those angels, the deserters from
God, the lovers of women,'' who ``were likewise the discoverers of this
curious art'' (\emph{De idol.}\ 9), and the Apostle's ``because of the
angels'' (1~Cor 11:10) means ``those, to wit, whom we read of as having
fallen from God and heaven on account of concupiscence after females''
(\emph{De virginibus velandis} 7; see \ref{sec:fathers-latin}). Cyprian
repeats the teaching of the arts: ``All which things sinning and apostate
angels put forth by their arts, when \dots\ they forsook their heavenly
vigour. They taught them also to paint the eyes with blackness drawn round
them in a circle'' (\emph{De habitu virginum} 14).

Lactantius gives the fullest Latin form, with the good beginning that
Jubilees also tells (\ref{sec:bc-jubilees}). God \latin{misit angelos ad
tutelam cultumque generis humani}, ``sent angels for the protection and
improvement of the human race; and inasmuch as He had given these a free
will, He enjoined them above all things not to defile themselves with
contamination from the earth.'' The devil enticed them by women, and their
offspring, ``neither angels nor men,'' were admitted neither to heaven nor to
hell: ``Thus there came to be two kinds of demons; one of heaven, the other
of the earth'' (\emph{Div.\ Inst.}\ II.14, CSEL; II.15 in ANF). Sulpicius
Severus still tells it in his \emph{Chronica} (I.2).

\paragraph{The Fathers who weighed the book.} Origen cites Enoch and
hesitates over him. Jared was born ``as it is written in the Book of
Enoch---if any one cares to accept that book as sacred---in the days when
the sons of God came down to the daughters of men'' (\emph{Comm.\ in Io.}\
VI.25; 1~En 6:6); against Celsus he notes that ``the books which bear the
name Enoch do not at all circulate in the Churches as divine'' (\emph{C.\
Cels.}\ V.54). Hilary knows the oath on Hermon and the meaning Enoch gives
the name. \emph{Hermon} is interpreted \latin{anathema}, and \latin{fertur
autem id, de quo etiam nescio cuius liber extat, quod angeli concupiscentes
filias hominum, cum de caelo descenderent, in hunc montem maxime excelsum
conuenerint}: it is related, and some book exists about it, that the angels
who desired the daughters of men came together on this highest mountain when
they descended from heaven. He goes no further: \latin{sed haec
praetermittamus. quae enim libro legis non continentur, ea nec nosse
debemus} (\emph{Tract.\ in Ps.}\ 132.6). Jerome records that Jude's letter
``because in it he quotes from the apocryphal book of Enoch it is rejected by
many. Nevertheless by age and use it has gained authority and is reckoned
among the Holy Scriptures'' (\emph{De viris illustribus} 4).

\paragraph{The Fathers who set the Watchers aside.} From the third century
the reading of Gen 6 as a history of angels gives way to the reading of the
sons of God as the line of Seth. Julius Africanus holds that ``the descendants of Seth are called the sons of
God on account of the righteous men and patriarchs who have sprung from
him''; if angels are meant, ``we must take them to be those who deal with
magic and jugglery, who taught the women the motions of the stars''
(\emph{Chronographia}, fragment 2).
Chrysostom sets out to overturn what he calls the tales of those who speak
without examination. He challenges them to show where angels are called
sons of God: Scripture, he says, gives the name to men (Ps 81[82]:6; Isa
1:2; Exod 4:22), never to angels. If they say the angels fell by this very
act, Scripture teaches that the
devil and his host fell before the first man was formed, since ``by the envy
of the devil, death came into the world'' (Wis 2:24); and the bodiless nature
could not receive such desire, for in the resurrection ``they shall neither
marry nor be married, but shall be as the angels of God in heaven'' (Matt
22:30). The sons of God are
those descended from Seth and from Enos, of whom it is written that he
``began to call upon the name of the Lord'' (Gen 4:26); the daughters of men
are those of Cain's line (\emph{Hom.\ in Gen.}\ 22.2--3).

Augustine gives the Latin Church its settled reading: ``certainly I could by
no means believe that God's holy angels could at that time have so fallen,''
and the angels that sinned of 2~Pet 2:4 are rather ``these who first
apostatized from God, along with their chief the devil'' (\emph{De civ.\
Dei} XV.23). Of the book he says: \latin{Scripsisse quidem nonnulla divine
illum Enoch, septimum ab Adam, negare non possumus, cum hoc in epistula
canonica Iudas apostolus dicat}, ``We cannot deny that Enoch, the seventh
from Adam, left some divine writings, for this is asserted by the Apostle
Jude in his canonical epistle''; but ``the writings which are produced under
his name, and which contain these fables about the giants, saying that their
fathers were not men, are properly judged by prudent men to be not
genuine'' (\emph{De civ.\ Dei} XV.23). They are not rejected ``because the
authority of these men who pleased God is rejected, but because they are not
believed to be theirs'' (XVIII.38). Cassian's Serenus holds that ``spiritual
existences'' cannot ``have carnal intercourse with women,'' and reads the
sons of God as Seth's holy line, called ``\,`angels of God,' or as some copies
have it `sons of God'\,''; the corrupt arts came through the daughters of
Cain and were carried past the Flood by Ham, ``In this way then that common
notion, according to which men believe that angels delivered to men
enchantments and diverse arts, is in truth fulfilled'' (\emph{Conl.}\
VIII.21). The arts are still the devils' teaching, taught through men.

\paragraph{Aquinas and the Church.} Aquinas answers the objection from Gen
6:4 in Augustine's words, ``God's holy angels could not fall in such fashion
before the deluge. Hence by the sons of God are to be understood the sons of
Seth'' (\ST{I q.~51 a.~3 ad~6}; see \ref{sec:q51}). The Church confesses
that \latin{Diabolus enim et alii daemones a Deo quidem natura creati sunt
boni, sed ipsi per se facti sunt mali} (Lateran IV, DS~800): the demons are
angels fallen by their own will, not wicked by nature (\ST{I q.~63 a.~4}),
and their sin was of the kind that belongs to a spirit, for ``a spiritual
nature cannot be affected by such pleasures as appertain to bodies'' (\ST{I
q.~63 a.~2}; see \ref{sec:fall}). The Watchers' history supplied
the Fathers with much that remains. The fallen are held in darkness until
the judgment (1~En 10:4--6, 12; Jude 6; 2~Pet 2:4), and Aquinas teaches that
the darksome atmosphere is their prison ``until the judgment day,'' while
``some of them are even now in hell'' (\ST{I q.~64 a.~4}, with Augustine; see
\ref{sec:q64}). The demons draw men to worship them as gods (1~En 19:1;
99:7), for ``all the gods of the Gentiles are devils'' (Ps 95[96]:5; 1~Cor
10:20). The holy angels carry the prayers of men (1~En 9:3; 40:6), as Raphael
does (Tob 12:12) and the angel with the golden censer (Apoc 8:3--4). The
righteous have guardians from among the holy angels (1~En 100:5), as the
little ones have angels who ``always see the face of my Father'' (Matt
18:10; \ST{I q.~113 a.~2}; see \ref{sec:q113}). The unions with women, the
giants, and the demons born of them the Fathers after Augustine set aside.

\sectionguard
\subsection{Jubilees, the Little Genesis}\label{sec:bc-jubilees}

The Book of Jubilees retells Genesis and the beginning of Exodus as a
revelation given to Moses on Sinai by ``the angel of the presence,'' who
dictates the history ``from the beginning of creation'' (Jub 1:27; 2:1). It
reached the Fathers in Greek as the \emph{Little Genesis}.\footnote{English
from R. H. Charles, \emph{The Book of Jubilees or the Little Genesis}
(London: SPCK, 1917), translated from the Ethiopic.} Epiphanius takes
from it the names of the wives of Cain and Seth, naming the book \gk{ta
I\=ob\=elaia}, which he says is also called the Little Genesis
(\emph{Panarion} 39.6.1, PG~41). Jerome, glossing a word of Numbers, finds it
nowhere in the Hebrew Scriptures \latin{absque libro apocrypho, qui a Graecis
\dots\ id est parva Genesis, appellatur} (\emph{Ep.}\ 78, eighteenth station,
PL~22).

\paragraph{The angels on the first day (Jub 2:2--3).} Jubilees gives the
angels' creation a day: ``For on the first day He created the heavens which
are above and the earth and the waters and all the spirits which serve
before Him---the angels of the presence, and the angels of sanctification,
and the angels [of the spirit of fire and the angels] of the spirit of the
winds, and the angels of the spirit of the clouds, and of darkness, and of
snow and of hail and of hoar frost, and the angels of the voices and of the
thunder and of the lightning, and the angels of the spirits of cold and of
heat, and of winter and of spring and of autumn and of summer, and of all
the spirits of His creatures which are in the heavens and on the earth
\dots'' (Jub 2:2). The angels are witnesses of the rest of the work: ``And thereupon we saw His
works, and praised Him, and lauded before Him on account of all His works;
for seven great works did He create on the first day'' (Jub 2:3), the image of the
Septuagint's Job (Job 38:7). Epiphanius carries the list into Christian
teaching without naming the book: the holy measure is the twenty-two works of
the six days, and on the first God made the upper heavens, the earth, the
waters, and \gk{ta pneumata ta leitourgounta en\=opion autou}, the spirits
that minister before him, angels of the presence and of glory, of winds and
clouds, of snow and hail, of thunders and lightnings, of cold, heat, and the
seasons, seven greatest works in all (\emph{De mensuris et ponderibus} 22,
PG~43; Jub 2:2--3, 23). Augustine reaches the first day from Scripture
alone (\emph{De civ.\ Dei} XI.9), and Aquinas holds as more probable that the
angels were created with the corporeal world (\ST{I q.~61 a.~3}).

The angels of wind, cloud, and season answer to the angels whom the Fathers
set over the elements (Athenagoras, \emph{Leg.}\ 24; see
\ref{sec:fathers-before}); Aquinas teaches that as the lower angels are
ruled by the higher, ``so are all corporeal things ruled by the angels,'' and
adds, ``This is not only laid down by the holy doctors, but also by all
philosophers who admit the existence of incorporeal substances'' (\ST{I
q.~110 a.~1}; see \ref{sec:government}). The angel of the presence ``who
went before the camp of Israel'' dictates the Law's history to Moses (Jub
1:29; Exod 14:19), in agreement with the Greek of Deut 33:2 and with the New
Testament that the Law came through angels (\ST{I-II q.~98 a.~3}).

\paragraph{The Watchers sent for good (Jub 4:15, 22; 5:1--10).} Jubilees
places the Watchers' coming in the days of Jared, like Enoch, but gives it a
good beginning: ``in his days the angels of the Lord descended on the earth,
those who are named the Watchers, that they should instruct the children of
men, and that they should do judgment and uprightness on the earth'' (4:15).
Enoch ``was moreover with the angels of God these six jubilees of years,''
and ``he testified to the Watchers, who had sinned with the daughters of
men'' (4:21--22). Against the angels ``whom He had sent upon the earth''
God commanded ``to root them out of all their dominion, and He bade us to
bind them in the depths of the earth'' until ``the day of the great
condemnation'' (5:6, 10). Justin and Lactantius give the same order of events, the angels first
appointed to the care of men and fallen in that charge (\emph{2~Apol.}\ 5;
\emph{Div.\ Inst.}\ II.14), and Irenaeus's Enoch testifies to the angels as
God's legate (\emph{Adv.\ haer.}\ IV.16.2); these are parallels, for none of
them names Jubilees.

\paragraph{Mastema and the tenth part (Jub 10:1--11; 17:16; 48).} After the
Flood the unclean demons lead Noah's grandchildren astray, and Noah prays
that they be bound, ``for Thou alone canst exercise dominion over them''
(10:6). ``The chief of the spirits, Mastêmâ,'' asks that some remain, ``for if some of them are not left to me, I shall not be
able to execute the power of my will on the sons of men''; God answers, ``Let
the tenth part of them remain before him, and let nine parts descend into
the place of condemnation'' (10:8--9). ``The prince Mastêmâ'' proposes the
testing of Abraham, as Satan asks for Job (17:16; Job 1:9--12), and he
strengthens Pharaoh's sorcerers, but the angel says, ``The evils indeed we
permitted them to work, but the remedies we did not allow to be wrought by
their hands'' (48:10). The book teaches that the power of the evil spirits
is permitted, limited, and ordered by God. Aquinas teaches the same of the
demons' assault: ``The assault itself is due to the malice of the demons
\dots\ But the ordering of the assault is from God, Who knows how to make
orderly use of evil by ordering it to good'' (\ST{I q.~114 a.~1}); and of
their place, that until the judgment ``the demons are in this dark
atmosphere for our trial: although some of them are even now in hell''
(\ST{I q.~64 a.~4}). The division between the spirits bound and those left
for man's trial is a parallel: Aquinas draws it from Augustine and the order
of providence.

\paragraph{The nations and their spirits (Jub 15:31--32).} Jubilees reads
the division of the nations otherwise than the Septuagint: ``over all hath
He placed spirits in authority to lead them astray from Him. But over Israel
He did not appoint any angel or spirit, for He alone is their ruler, and He
will preserve them and require them at the hand of His angels'' (15:31--32).
The Fathers who read Deut 32:8 in the Greek hold the contrary on both
points. Dionysius denies that the nations' wandering is to be laid ``upon
the direct guidance of the Angels'' and teaches that the angels of the
nations lead their peoples toward the one God (\emph{CH} 9.3--4); Scripture
and Dionysius give Israel an angelic prince, Michael (Dan 10:21; 12:1;
\emph{CH} 9.2); Eusebius gives Israel to the Word himself (\emph{Dem.\
evang.}\ IV.7). The Catechism teaches with the Greek Song that the nations are entrusted
``to the guardianship of angels'' (CCC~57), and Aquinas that their princes
are good angels executing God's judgments (\ST{I q.~113 a.~8}).

\sectionguard
\subsection{The Testaments of the Twelve Patriarchs}\label{sec:bc-testaments}

The \emph{Testaments of the Twelve Patriarchs} give the dying counsels of
the twelve sons of Jacob. The book reached the Church in a Greek text that
names Christ in many places; Sinker, its translator for the Ante-Nicene
library, judged its author a Jew converted to Christianity.\footnote{English
from R. Sinker's translation, \emph{Ante-Nicene Fathers} vol.~8 (American
edition, 1886), cited by testament and chapter.} Origen cites it by name:
\latin{Sed et in aliquo quodam libello, qui appellatur testamentum duodecim
patriarcharum, quamvis non habeatur in canone, talem tamen quendam invenimus
sensum, quod per singulos peccantes singuli Satanae intelligi debeant}; in a
little book called the Testament of the Twelve Patriarchs, though it is not
in the canon, we find the sense that for each sinner a satan is to be
understood (\emph{Hom.\ in Iesu Nave} 15.6, GCS 7).

\paragraph{The heavens and their ranks (T.~Levi 2--5).} Levi is taken up
through seven heavens. The lower hold the instruments of judgment, the
higher are holy: ``In the heaven next to it are the angels of the presence
of the Lord, who minister and make propitiation to the Lord for all the
ignorances of the righteous; and they offer to the Lord a reasonable
sweet-smelling savour, and a bloodless offering \dots\ And in the heaven next
to this are thrones, dominions, in which hymns are ever offered to God''
(T.~Levi 3). The angels are ranked,
the highest nearest God, their service is worship, and among the ranks are
the Apostle's thrones and dominions (Col 1:16). The offering of the angels of
the presence answers to the angel of the Apocalypse who offers ``the prayers
of all saints'' (Apoc 8:3). Origen gives the same service in Christian terms: the
angels ``ascend, bearing the supplications of men, to the purest of the
heavenly places in the universe,'' and descend with God's gifts; yet ``every
prayer, and supplication, and intercession, and thanksgiving, is to be sent
up to the Supreme God through the High Priest, who is above all the angels,
the living Word and God'' (\emph{C.\ Cels.}\ V.4).

\paragraph{The angel who intercedes (T.~Levi 5; T.~Dan 6).} The angel who
leads Levi names himself: ``I am the angel who intercedeth for the race of
Israel, that He smite them not utterly, because every evil spirit attacketh
it'' (T.~Levi 5). Dan bids his sons ``draw near unto God, and to the Angel
that intercedeth for you, for He is a Mediator between God and man for the
peace of Israel. He shall stand up against the kingdom of the enemy'' (T.~Dan
6). The office of the intercessor is Michael's in Daniel (Dan 12:1) and in
Hermas (\emph{Sim.}\ VIII.3; see \ref{sec:fathers-before}). The name of
Mediator the Church gives to Christ alone: ``there is one God: and one
mediator of God and men, the man Christ Jesus'' (1~Tim 2:5). Aquinas makes
the verse the \emph{sed contra} that mediation is proper to Christ, and
grants the angels a share in it only as ministers: they ``fulfill the office
of mediator, not indeed principally and perfectively, but ministerially and
dispositively'' (\ST{III q.~26 a.~1} and ad~2, where he cites Augustine,
\emph{De civ.\ Dei} IX.13). The Testaments' intercessor is received as the
ministering angel of Israel and corrected where the name Mediator would
place him beside the Son.

\paragraph{The two spirits and the spirits of error.} ``Two spirits wait upon
man---the spirit of truth and the spirit of error; and in the midst is the
spirit of the understanding of the mind, to which it belongeth to turn
whithersoever it will'' (T.~Judah 20). A man's end shows whether he knew
``the angels of the Lord from the angels of Satan,'' for the soul that served
evil ``is tormented by the evil spirit,'' but if ``it hath known the angel of
peace, it shall comfort him in life'' (T.~Asher 6). ``The
mind of the good man is not in the power of the deceit of the spirit of
Beliar, for the angel of peace guideth his soul'' (T.~Benjamin 6). Reuben
counts ``seven spirits of error,'' first ``the spirit of fornication''
(T.~Reuben 2--3). The Fathers carried these teachings. Hermas taught that ``There are
two angels with a man---one of righteousness, and the other of iniquity''
(\emph{Mand.}\ VI.2), and Cassian's Serenus, citing Hermas with Matt 18:10,
teaches that ``two angels, a good and a bad one, cling to each one of us''
(\emph{Conl.}\ VIII.17). Origen, in the homily that names the Testaments,
teaches that the adverse spirits are ranked by the vices: \latin{est aliqui
fornicationis spiritus, est et irae; spiritus alius est avaritiae, alius vero
superbiae}, a spirit of fornication, of anger, of avarice, of pride; each
kind has one prince, \latin{innumeros vero esse, qui in hoc ei officio
pareant}, and innumerable spirits who obey him in that office; and above them
all stands one prince of all wickedness (\emph{Hom.\ in Iesu Nave} 15.5).
Aquinas receives the order of the assault and corrects its extent. The
demons ``through pride usurp a semblance of Divine power, by deputing
certain ministers to assail man, as the angels of God in their various
offices minister to man's salvation'' (\ST{I q.~114 a.~1}); yet the devil is
the direct cause of some sins only, ``for all sins are not committed at the
devil's instigation, but some are due to the free-will and the corruption of
the flesh,'' as Origen himself teaches (\ST{I q.~114 a.~3}, citing \emph{Peri
Archon} III; see \ref{sec:q114}). The good angel who guides the soul is the
guardian angel of the Church's teaching (\ST{I q.~113 aa.~1--2}).

The Testaments also keep the Watchers. The women ``allured the Watchers
before the flood,'' who ``changed themselves into the shape of men'' (T.~Reuben
5), and ``the Watchers changed the order of their nature, whom also the Lord
cursed at the flood'' (T.~Naphtali 3; see \ref{sec:bc-enoch}).

\sectionguard
\subsection{The Seven Who Stand before God}\label{sec:bc-seven}

Tobias and the Books of the Machabees are Scripture for the Church, and
their angels are expounded with the canon (\ref{sec:scripture}). One
teaching of Tobias meets the Enochic books. Raphael names his rank: ``For I
am the angel Raphael, one of the seven, who stand before the Lord'' (Tob
12:15). The Greek of Tobias adds their office: the seven holy angels present the
prayers of the saints and go in and out before the glory of the Holy One (Tob
12:15, Greek). John sees ``seven angels standing in the presence
of God'' (Apoc 8:2). Enoch names seven ``holy angels who watch'': Uriel,
``over the world and over Tartarus''; Raphael, ``over the spirits of men'';
Raguel; Michael, ``set over the best part of mankind''; Saraqâêl; Gabriel,
``over Paradise and the serpents and the Cherubim''; and Remiel (1~En
20:1--8). Clement of Alexandria teaches that ``The first-born princes of the
angels, who have the greatest power, are seven'' (\emph{Strom.}\ VI.16).
Origen ascribes to three of them the offices Enoch gives: ``to Raphael \dots\
the work of curing and healing; to Gabriel, the conduct of wars; to Michael,
the duty of attending to the prayers and supplications of mortals''
(\emph{De princ.}\ I.8.1; 1~En 40:9), though he grounds the offices in
merit earned before the world (see \ref{sec:fathers-before}). Gregory holds
that the angels take their names from their ministries when they come to us,
\latin{apud nos etiam nomina a ministeriis trahunt} (\emph{Hom.\ in Ev.}\
34.8--9; see \ref{sec:scripture}).

The names Enoch adds did not enter the Church's prayer. In 745 the Roman
synod under Pope Zachary heard read the prayer of Aldebert, which invoked eight
angels by name, beginning with Uriel and Raguel; the bishops answered that,
Michael excepted, Aldebert had invoked demons under the names of angels, and
that
\latin{non plus quam trium angelorum nomina cognoscimus, id est Michael,
Gabriel, Raphael} (Concilium Romanum, a.~745, MGH \emph{Concilia} II.1, pp.~42--43).
The \emph{Directory on Popular Piety and the Liturgy} (2001) repeats the
rule: ``The practice of assigning names to the Holy Angels should be
discouraged, except in the cases of Gabriel, Raphael and Michael whose names
are contained in Holy Scripture'' (no.~217). The seven are Scripture's; of their
names the Church keeps three (see \ref{app:names}; \ref{sec:devotion}).

\sectionguard
\subsection{Philo of Alexandria}\label{sec:bc-philo}

Philo, the Jewish philosopher of Alexandria and a contemporary of the
Apostles, expounded the Law of Moses in the language of the Greek schools,
and the Fathers read him. Clement names ``the Pythagorean Philo'' as a
witness to the antiquity of the Hebrews' philosophy (\emph{Strom.}\ I.15);
Origen holds that the allegorists of the Law, Philo among them, ``have so
successfully hit the meaning (of the sacred writers), that even Grecian
philosophers would have been captivated'' (\emph{C.\ Cels.}\ IV.51); Jerome
places him ``among the ecclesiastical writers'' and lists among his works
``Concerning the confusion of tongues,'' ``Concerning giants,'' and ``That
dreams are sent by God'' (\emph{De vir.\ ill.}\ 11).\footnote{English of
Philo from C. D. Yonge, \emph{The Works of Philo Judaeus} (London: Bohn,
1854--1855), cited by the section numbers of the modern editions.} These are
the treatises in which he speaks of the angels.

\paragraph{Angels, daemons, and souls.} On Gen 6:2, which he reads ``the
angels of God,'' Philo writes: ``Those beings, whom other philosophers call
demons, Moses usually calls angels; and they are souls hovering in the air''
(\emph{De gigantibus} 6). Some souls descend into bodies; others the Creator
employs ``as hand-maidens and servants in the administration of mortal
affairs'' (12), the good among them ``ambassadors of man to God, and of God
to man, and sacred and holy on account of this blameless and most excellent
office'' (16). On Jacob's ladder he says of the purest souls: ``philosophers
in general are wont to call these demons, but the sacred scripture calls
them angels, using a name more in accordance with nature. For indeed they do
report (\gk{diangellousi}) the injunctions of the father to his children,
and the necessities of the children to the father.'' They ascend and
descend ``not because God \dots\ has any need of interpreters; but because it
is the lot of us miserable mortals to use speech as a mediator and
intercessor'' (\emph{De somniis} I.141--142; cf.\ \emph{De plantatione} 14).

The Fathers received the teaching that the angels are named by their
office. Origen: ``Having thus learned to call these beings `angels' from their
employments, we find that because they are divine they are sometimes termed
`god' in the sacred Scriptures, but not so that we are commanded to honour
and worship in place of God those who minister to us'' (\emph{C.\ Cels.}\
V.4). Augustine gives it its classic form: \latin{Angelus enim officii nomen
est, non naturae} (\emph{Enarr.\ in Ps.}\ 103, s.~1, 15; see
\ref{sec:scripture}). Philo's making of angels, daemons, and souls one kind
the Church rejects. Aquinas reports that Origen ``held that human souls and
angels are all of the same species'' (\ST{I q.~75 a.~7}, citing \emph{Peri
Archon} III.5); the canons against Origen issued at Constantinople in 543
anathematize whoever says \latin{praeexsistere hominum animas, utpote quae
antea mentes fuerint et sanctae virtutes}, that men's souls pre-existed as
minds and holy powers and were sent into bodies for punishment (DS~403); and
Aquinas determines that soul and angel differ in species, since in
incorporeal substances ``there cannot be diversity of number without
diversity of species'' (\ST{I q.~75 a.~7}). Origen remembers Philo's reading
of Gen 6 without naming him: ``even before us there was one who referred this
narrative to the doctrine regarding souls, which became possessed with a
desire for the corporeal life of men'' (\emph{C.\ Cels.}\ V.55).

\paragraph{The king and his powers.} Philo explains the plurals of ``Let us
make man'' and ``Come, let us go down and confuse'' by God's powers. God is
one, but ``has about him an unspeakable number of powers''; in the air is ``a
most sacred company of incorporeal souls \dots\ for the word of prophecy is
accustomed to call these souls angels,'' marshalled ``in their suitable
ranks'' as ``servants and ministers of the ruler'' (\emph{De confusione
linguarum} 171--174). ``It is suitable to the character of the king to
associate with his own powers,'' though ``the Father of the universe has no
need of anything'' (175). Aquinas gives the ministry of the angels a deeper
reason. God governs all things immediately in the design of his providence,
but in its execution ``He governs some things by means of others,'' because
``it is a greater perfection for a thing to be good in itself and also the
cause of goodness in others''; and ``That an earthly king should have
ministers to execute his laws is a sign not only of his being imperfect, but
also of his dignity'' (\ST{I q.~103 a.~6} and ad~3; see
\ref{sec:government}). Philo grounds the ministry in God's dignity and man's
weakness, Aquinas in the communication of God's goodness.

\paragraph{Whether the angels made man.} Philo goes further: ``Very
appropriately therefore has God attributed the creation of this being, man,
to his lieutenants, saying, `Let us make man,' in order that the successes
of the intellect may be attributed to him alone, but the errors of the being
thus created, to his subordinate power'' (\emph{De conf.}\ 179; cf.\ \emph{De
opificio mundi} 72--75). The Fathers reject it. Justin will not grant
``that the dogma of that heresy which is said to be among you is true, or
that the teachers of it can prove that [God] spoke to angels, or that the
human frame was the workmanship of angels''; the Father spoke to the Son
(\emph{Dial.}\ 62). Irenaeus: ``It was not angels, therefore, who made us,
nor who formed us, neither had angels power to make an image of God \dots\
For with Him were always present the Word and Wisdom, the Son and the Spirit''
(\emph{Adv.\ haer.}\ IV.20.1). Aquinas determines that creation belongs to
God alone (\ST{I q.~45 a.~5}) and that ``We must not imagine that when God
said `Let us make man,' He spoke to the angels, as some were perverse enough
to think'' (\ST{I q.~91 a.~4 ad~2}); the angels could ``act as
ministers in the formation of the body of the first man, in the same way as
they will do at the last resurrection by collecting the dust'' (\ST{I q.~91
a.~2 ad~1}).

\paragraph{The giants.} In the \emph{Questions on Genesis} Philo reads Gen
6:4 of angels: the giants ``were sprung from a combined procreation of two
natures, namely, from angels and mortal women,'' and ``sometimes Moses styles
the angels the sons of God, inasmuch as they were not produced by any
mortal'' (\emph{Quaest.\ in Gen.}\ I.92). Ambrose follows him: \latin{non
poetarum more gigantas illos terrae filios uult uideri diuinae scripturae
conditor, sed ex angelis et mulieribus generatos adserit}, Scripture means
not the poets' sons of the earth but men begotten of angels and women
(\emph{De Noe} 4.8, CSEL 32.1, where Schenkl marks the source in Philo). The
Latin Church after Augustine sets the reading aside (\ref{sec:bc-enoch}).

\sectionguard
\subsection{Israel's Writings in the Fathers}\label{sec:bc-israel-reception}

From Israel's writings outside the canon the Fathers kept the substance of
the angels' ministry: the angels set over peoples and over the elements,
their praise at the founding of the world, the Law given through them, their
intercession and the carrying of prayers, the guardian of the just, the
seven who stand before God, and the binding of the fallen in darkness until
the judgment. They refused the marriages of angels with women, the demons as
spirits of giants, an angel called mediator beside Christ, and man made by
angels. The Septuagint's angels they received as Scripture.

\begin{fourstable}{Writing and locus}{Teaching}{Patristic reception}{Disposition}
Deut 32:8--9, Septuagint & The nations bounded by the number of the angels of God; Israel the Lord's portion & \emph{1 Clem.}\ 29; Justin, \emph{Dial.}\ 131; Clement, \emph{Strom.}\ VII.2; Origen, \emph{De princ.}\ I.5.2; Basil, \emph{Adv.\ Eun.}\ III.1; Eusebius, \emph{Dem.\ ev.}\ IV.7; Dionysius, \emph{CH} 9.2--4 & Received: CCC~57; \ST{I q.~108 a.~6}; q.~113 a.~8; in Christ the nations pass to the Lord (Irenaeus III.12.9; Hilary, \emph{Tr.\ Ps.}\ 2.31)\\
Job 1:6; 2:1, Septuagint & The angels of God stand before the Lord, the devil among them & Gregory, as Aquinas cites him & Received: \ST{I q.~112 a.~3 ad~3}\\
Job 38:7, Septuagint & All the angels praised God when the stars were made & Augustine, \emph{De civ.\ Dei} XI.9 (first day); Cassian, \emph{Conl.}\ VIII.7 (before the world) & Received; created with the corporeal world, the common teaching (\ST{I q.~61 a.~3})\\
Deut 32:43; Ps 8:6; 77[78]:25; 96[97]:7; 137[138]:1, Septuagint and Vulgate & The angels adore the Son; man a little less than the angels; bread of angels; psalmody before the angels & Heb 1:6; 2:7; Wis 16:20; Benedict, \emph{Regula} 19 & Received as Scripture\\
Deut 33:2, Septuagint; Jub 1:27--2:1 & The angels at Sinai; the angel of the presence gives the Law & Acts 7:53; Gal 3:19; Heb 2:2 & Received: \ST{I-II q.~98 a.~3}\\
Isa 9:6; 63:9, Septuagint & The Angel of great counsel; not an angel but the Lord himself saves & Justin, \emph{Dial.}\ 126; Origen, \emph{C.\ Cels.}\ V.53; Irenaeus III.20.4 & Received of Christ\\
Dan 3:49, 58; 4:10, 14 & The angel in the furnace; the angels bless the Lord; the watchers & Augustine, \emph{De civ.\ Dei} XI.9; Jerome, \emph{In Dan.}\ 4:10 & Received as Scripture\\
1~En 1:9 & The Lord comes with his holy ones to judge & Jude 14--15; Tertullian, \emph{Cult.\ fem.}\ I.3; Augustine, \emph{De civ.\ Dei} XV.23 & Received through Jude; book Apocryphal (Trent, DS~1503; Jerome, \emph{Vir.\ ill.}\ 4)\\
1~En 6--8; 15--16 & Watchers took wives, begot giants, taught arts; giants' spirits are demons & Justin, \emph{2~Apol.}\ 5; Athenagoras, \emph{Leg.}\ 24--25; Irenaeus IV.16.2; Clement, \emph{Strom.}\ V.1; Tertullian, \emph{Cult.\ fem.}\ I.2; Cyprian, \emph{Hab.\ virg.}\ 14; Lactantius II.14 & Rejected: Africanus, frag.\ 2; Chrysostom, \emph{Hom.\ in Gen.}\ 22; Augustine XV.23; Cassian VIII.21; \ST{I q.~51 a.~3 ad~6}; q.~63 a.~2; Lateran IV (DS~800)\\
1~En 10:4--12; Jub 5:6--10 & The fallen bound in darkness until the judgment & Jude 6; 2~Pet 2:4; Augustine XV.23 & Received: \ST{I q.~64 a.~4}\\
1~En 19:1; 99:7 & Fallen spirits lead men to sacrifice to demons as gods & Tertullian, \emph{Idol.}\ 4 (quoted as Enoch) & Received with Ps 95[96]:5; 1~Cor 10:20\\
1~En 9:3; 40:6 & Archangels carry men's cause and prayers before God & Origen, \emph{C.\ Cels.}\ V.4 & Received with Tob 12:12; Apoc 8:3--4\\
1~En 14:22--23; 39:12; 61:10; 71:7 & Those who sleep not sing the Sanctus; angels of powers and principalities & Jerome, \emph{In Dan.}\ 4:10 (watchers) & Received with Isa 6:3; Dan 7:10; Col 1:16 (\ST{I q.~108 a.~5})\\
1~En 100:5 & Guardians from the holy angels for the righteous & Parallel only & Received with Matt 18:10 (\ST{I q.~113 a.~2})\\
1~En 20; 40:9 & Seven named archangels; Raphael over diseases & Clement, \emph{Strom.}\ VI.16; Origen, \emph{De princ.}\ I.8.1 & Seven received (Tob 12:15; Apoc 8:2); names beyond three rejected (Roman synod 745; \emph{Directory} 217)\\
Jub 2:2--3 & The angels created on the first day, praising God's works & Epiphanius, \emph{De mens.}\ 22; Augustine, \emph{De civ.\ Dei} XI.9 (parallel) & Created with the corporeal world, the common teaching (\ST{I q.~61 a.~3}); angels over the elements received (\ST{I q.~110 a.~1})\\
Jub 4:15 & Watchers first sent to instruct men & Justin, \emph{2~Apol.}\ 5; Lactantius II.14 (parallel) & Corrected: good angels are sent to men (\ST{I q.~112}); the unions rejected\\
Jub 10:8--11; 48:10 & A tenth of the spirits left to Mastema; evils permitted and limited & Parallel only & Received in substance: \ST{I q.~114 a.~1}; q.~64 a.~4\\
Jub 15:31--32 & Spirits over the nations lead them astray; no angel over Israel & Dionysius, \emph{CH} 9.3--4; Eusebius, \emph{Dem.\ ev.}\ IV.7 & Corrected: CCC~57; \ST{I q.~113 a.~8}\\
T.~Levi 3 & Ranks in the heavens; angels of the presence offer a bloodless offering; thrones, dominions & Origen, \emph{C.\ Cels.}\ V.4 & Received in substance (Apoc 8:3; Col 1:16)\\
T.~Levi 5; T.~Dan 6 & The angel who intercedes for Israel, called Mediator & Hermas, \emph{Sim.}\ VIII.3 (Michael) & Corrected: one Mediator (1~Tim 2:5); angels mediate ministerially (\ST{III q.~26 a.~1 ad~2})\\
T.~Judah 20; T.~Asher 6; T.~Benj.\ 6 & Two spirits; angels of the Lord and of Satan; the angel of peace guides the soul & Hermas, \emph{Mand.}\ VI.2; Cassian, \emph{Conl.}\ VIII.17 & Received: \ST{I q.~113 aa.~1--2}; q.~114 a.~1\\
T.~Reuben 2--3 & Seven spirits of error, beginning with fornication & Origen, \emph{Hom.\ in Iesu Nave} 15.5--6 (names the book) & Corrected: \ST{I q.~114 a.~3}\\
Philo, \emph{De gig.}\ 16; \emph{De somn.}\ I.141--142 & Angels named for their office as messengers & Origen, \emph{C.\ Cels.}\ V.4; Augustine, \emph{Enarr.\ in Ps.}\ 103 & Received\\
Philo, \emph{De gig.}\ 6--16 & Angels, daemons, and souls one kind & Origen, \emph{C.\ Cels.}\ V.55 & Rejected: DS~403; \ST{I q.~75 a.~7}\\
Philo, \emph{De conf.}\ 171--175 & God rules through his powers as a king through ministers & Parallel & Received and deepened: \ST{I q.~103 a.~6}\\
Philo, \emph{De conf.}\ 179; \emph{De opif.}\ 72--75 & ``Let us make man'' said to subordinate powers & Justin, \emph{Dial.}\ 62; Irenaeus IV.20.1 & Rejected: \ST{I q.~45 a.~5}; q.~91 a.~4 ad~2\\
Philo, \emph{Quaest.\ in Gen.}\ I.92 & Giants from angels and women & Ambrose, \emph{De Noe} 4.8 & Received by Ambrose; rejected after Augustine\\
\end{fourstable}
